三段代码，三种症状。一个序列模型把一千个词的条件概率连乘起来，似然返回 0，困惑度是 \texttt{inf}。一个自己写的反向传播，损失在降，但比参考实现慢一截，看不出哪里错。一个五百维高斯的对数似然，先求逆再求行列式，最后 \(\log\det\) 给出 \(-\infty\)。

三段代码都是照着公式一字一句翻译过来的，公式本身也都对。

在实数上，一个数学表达式只有一个值。在浮点数上，通往这个值的路径有许多条，它们的可用范围、算术代价和误差放大倍数差得很远——差到有的路径根本走不通。这一章练的就是识路：先辨认最终要的是什么，再挑一条不制造巨大中间量的走法。

\begin{center}
\begin{tikzpicture}[x=1cm,y=1cm,font=\small,>=stealth]
  \node[font=\footnotesize,black!70] at (2.15,3.75) {照公式直译};
  \node[font=\footnotesize,black!70] at (6.8,3.75) {症状};
  \node[font=\footnotesize,black!70] at (11.25,3.75) {本章的改写};
  \foreach \ya in {2.6,1.4,0.2}{
    \draw[fill=black!7,draw=black!45] (0,\ya) rectangle (4.3,\ya+0.9);
    \draw[fill=orange!16,draw=black!45] (5.2,\ya) rectangle (8.4,\ya+0.9);
    \draw[fill=blue!12,draw=black!45] (9.3,\ya) rectangle (13.2,\ya+0.9);
    \draw[->,black!60] (4.4,\ya+0.45) -- (5.15,\ya+0.45);
    \draw[->,black!60] (8.5,\ya+0.45) -- (9.25,\ya+0.45);
  }
  \node[font=\footnotesize,align=center] at (2.15,3.05) {\(\prod_t p_t\)，\(\log\sum_i e^{x_i}\)};
  \node[font=\footnotesize,align=center] at (6.8,3.05) {下溢成 0\\上溢成 NaN};
  \node[font=\footnotesize,align=center] at (11.25,3.05) {换到对数域\\先减最大值};
  \node[font=\footnotesize,align=center] at (2.15,1.85) {\(\dfrac{f(x+h)-f(x)}{h}\)};
  \node[font=\footnotesize,align=center] at (6.8,1.85) {截断与舍入两头夹\\还要跑 \(n\) 遍};
  \node[font=\footnotesize,align=center] at (11.25,1.85) {换到计算图\\逐节点链式法则};
  \node[font=\footnotesize,align=center] at (2.15,0.65) {\(A^{-1}b\)，\(\log\det A\)};
  \node[font=\footnotesize,align=center] at (6.8,0.65) {三倍算术，结构尽失\\行列式先溢出};
  \node[font=\footnotesize,align=center] at (11.25,0.65) {换成解方程\\一次分解取三样};
\end{tikzpicture}

\emph{三篇的分工。左列是数学符号的逐字翻译，右列是同一个量的另一种算法——两列在实数上完全相等，在浮点上一列能跑、一列不能。}
\end{center}

第一篇处理范围。概率连乘的下溢不是意外，是算术上的必然：一千个 \(10^{-8}\) 相乘等于 \(10^{-8000}\)，这个数 float64 从来就装不下。搬到对数域，连乘变连加，范围问题消失；可概率模型总要做边缘化，于是对数域里冒出 \(\log\sum_i e^{x_i}\)，直接算又会两头出界。log-sum-exp 提出最大值这一步，是稳定性改造最干净的样本：数学上是恒等式，数值上是能跑与不能跑之别。

第二篇处理导数。有限差分需要给每个参数跑一遍模型，还要在截断误差与舍入误差之间挑一个折中的步长；自动微分不走这条路，它对程序的计算图逐节点套链式法则，得到的导数精确到浮点运算本身的精度，那条 U 形误差曲线被整个绕开了。前向还是反向，由输入输出的维数决定——一个标量损失、上亿个参数，答案只能是反向。有限差分并没有出局，它退到了另一个岗位上：用一条独立的计算路径去检查求导程序有没有写错。

第三篇处理结构。\(x=A^{-1}b\) 是数学记号，不是操作指令。显式求逆要多花三倍算术、把稀疏或三角结构抹成稠密矩阵，得到的解连残差都不小。协方差的二次型应当由 Cholesky 因子做三角求解，\(\log\det\) 应当从同一个因子的对角线上加出来，而不是先算行列式再取对数——那一步会先溢出。

\begin{enumerate}
\tightlist
\item
  在对数域中稳定计算概率
\item
  自动微分给出导数，梯度检查建立信任
\item
  需要的是解，而不是逆矩阵
\end{enumerate}

三篇共用一条判断，它比任何一个具体技巧都值钱：看见一个公式，先问它是在陈述一个等式，还是在描述一条计算路径。绝大多数时候是前者，而把前者当成后者，就是本章三种症状的共同来源。

\subsection{本章篇目}

\begin{itemize}
\tightlist
\item \hyperref[sec:08-Numerical-Analysis/07-Numerical-Methods-for-ML/01]{``在对数域中稳定计算概率''}：连乘为什么必然下溢，log-sum-exp 怎样把可用窗口搬到输入身边，以及交叉熵为何不该绕道概率表；
\item \hyperref[sec:08-Numerical-Analysis/07-Numerical-Methods-for-ML/02]{``自动微分给出导数，梯度检查建立信任''}：计算图上的两种遍历方向，代价由维数决定，以及一次有判断力的梯度检查该怎么搭；
\item \hyperref[sec:08-Numerical-Analysis/07-Numerical-Methods-for-ML/03]{``需要的是解，而不是逆矩阵''}：求逆的三笔账，一次 Cholesky 服务三项任务，以及隐式层的反传为什么还是解方程。
\end{itemize}

需要更多背景的话，浮点数的表示与抵消现象在\hyperref[ch:061]{``浮点数与误差：计算机为什么不等于实数系统''}里；分解、主元与迭代解法的完整讨论在\hyperref[ch:065]{``数值线性代数：结构、稳定性与规模决定算法''}里。前一单元：\hyperref[ch:066]{``常微分方程数值方法：沿局部斜率走完整条轨迹''}。
